\section{MENTOR a porovnanie so SER}
\subsection{Ako funguje MENTOR}
MENTOR má tri hlavné úlohy. Najprv hľadá udalosti, o ktorých sa diskutuje v správach. Potom predikuje budúce udalosti. Nakoniec odhaduje poradie výkonnosti odvetvových indexov. Odvetvový index sleduje skupinu firiem z jedného odvetvia. Výstup teda hovorí, ktoré odvetvia majú podľa modelu dosiahnuť lepší výsledok. Nejde o percentuálny výnos jednej akcie \mentorref{kap.~3, rovnice~(1)--(3)}.

Článok používa aj pojem naratív. Ide o spôsob, akým sa udalosti opisujú a spájajú do spoločného príbehu alebo témy. Tento príbeh sa môže časom meniť. MENTOR používa viac agentov, teda častí systému s rôznymi úlohami. Agent na hľadanie udalostí spracúva veľa správ. Používa krátkodobú a dlhodobú pamäť. Vybrané udalosti potom pomocou vyhľadávania a zoradenia spája s pôvodnými správami. Autori tento postup porovnávajú aj so zoskupovaním podobných textových vektorov \mentorref{kap.~3--4.2}.

Ďalší agenti majú úlohy študenta, učiteľa a agenta pre uvažovanie. Študenti predikujú udalosti a poradie odvetví. Používajú textové stratégie $S_{\mathrm{event}}$ a $S_{\mathrm{ranking}}$. Sú to postupy opísané slovami, podľa ktorých vytvárajú predikciu. Keď sú známe skutočné výsledky, učitelia opíšu chyby. Agenti pre uvažovanie potom podľa tohto hodnotenia upravia textové stratégie \mentorref{kap.~4.4, rovnice~(4)--(7)}.

Tento postup využíva TextGrad. Spätná väzba má podobu textu a mení textové stratégie. Nejde o bežné spätné šírenie číselnej chyby, ktoré by menilo váhy LLM. Podľa algoritmu~1 sa výsledok týždňa $n$ predikuje z informácií predchádzajúceho týždňa. Skutočný výsledok týždňa $n$ sa použije až neskôr na úpravu stratégie. Pri predikcii teda ešte nesmie byť známy \mentorref{kap.~4.4, algoritmus~1}.

\begin{figure}[htbp]
\centering
\begin{tikzpicture}
\node[box,text width=4.3cm,minimum height=1.2cm] (in) at (0,0) {Správy a udalosti\\do týždňa $n-1$};
\node[box,text width=4.3cm,minimum height=1.2cm] (pred) at (5.5,0) {Študenti predikujú:\\udalosti a poradie};
\node[box,text width=4.3cm,minimum height=1.2cm] (feedback) at (5.5,-1.9) {Výsledok týždňa $n$:\\učitelia hodnotia};
\node[box,text width=4.3cm,minimum height=1.2cm] (update) at (0,-1.9) {Agenti pre uvažovanie:\\úprava stratégií};
\draw[arr] (in)--(pred);
\draw[arr] (pred)--(feedback);
\draw[arr] (feedback)--(update);
\draw[arr,dashed] (update.west)--++(-0.55,0)|- (in.west);
\end{tikzpicture}
\caption{Cyklus MENTOR: predikcia týždňa $n$ a úprava stratégií pre ďalšie obdobie. Vlastné spracovanie podľa \mentorref{obr.~3 a algoritmus~1}.}
\label{fig:mentor}
\end{figure}


\subsection{Porovnanie a výsledky}
\begin{table}[H]
\centering
\caption{Porovnanie oboch metód. Vlastné spracovanie podľa \serref{kap.~2--3} a \mentorref{kap.~3--5}.}
\label{tab:comparison}
\small
\begin{tabularx}{\textwidth}{@{}p{2.35cm}YY@{}}
\toprule
Vlastnosť & SER & MENTOR\\
\midrule
Vstup & Články spojené s firmou a identifikátor jej akcie & Správy, známe udalosti a minulé výnosy odvetví\\
Výstup & Odhad výnosu akcie na ďalší deň alebo týždeň & Budúce udalosti a poradie odvetvových indexov\\
Úloha LLM & Získanie udalostí z textu a ich kontrola & Hľadanie a predikcia udalostí, hodnotenie chýb\\
Čo sa mení & Vektory, attention a MLP podľa chyby výnosu & Textové postupy podľa hodnotenia učiteľov\\
Overenie & Test na ďalšom roku, výsledky portfólií a regresie & Testy po týždňoch, zhoda udalostí, poradie a portfóliá\\
\bottomrule
\end{tabularx}
\end{table}

Autori testujú MENTOR na čínskych príspevkoch vplyvných autorov a na anglických finančných správach z rokov 2023--2024. Pri predikcii udalostí zisťujú, ktoré predikované udalosti sa zhodujú so skutočnými. Pri poradí odvetví používajú Spearmanovu koreláciu, Kendallov koeficient a zhodu odvetví na prvých miestach. Prvé dve metriky porovnávajú predikované poradie so skutočným \mentorref{kap.~5.1--5.2}. Tieto testy však neoverujú, či by napríklad predikcia pravdepodobnosti 70\,\% zodpovedala približne 70\,\% skutočných prípadov. To by bola samostatná kontrola kalibrácie pravdepodobností.

Na 11 odvetviach S\&P dosahuje MENTOR\_O1 Spearmanovu koreláciu 0,221. Metóda Momentum dosahuje 0,159. Pri zhode troch najvyšších odvetví sú výsledky 0,394 a 0,364. Na deviatich odvetviach A-share má MENTOR\_O1 koreláciu 0,173. Pri zhode prvých troch odvetví však dosahuje 0,439 oproti 0,444 pri Momentum. MENTOR teda nie je lepší vo všetkých ukazovateľoch \mentorref{tab.~3}.

Autori uvádzajú aj výsledky portfólií a skúšajú rôzne dvojice modelov učiteľ--študent \mentorref{kap.~5.6, tab.~5--6}. Tieto výnosy nemožno priamo porovnať s výsledkami SER. Články používajú iné obdobia, aktíva a testy. Niektoré presné definície metrík a vzory pokynov sú v samostatných prílohách MENTOR. Pri príprave tejto práce tieto prílohy neboli dostupné na overenie.
